%%
%% Hypothesentest
%% 2023 - 05 - 24 φ
%%

\section{Testen von Hypothesen (optional)}\index{Testen von Hypothesen}\index{Hypothesentest}


Lernziele

\begin{itemize}
\item Nullhypothese $H_0$ = Alles ist zufällig bzw. gleich verteilt.

\item Alternativhypothese ($H_1$) = Es kann kein Zufall sein
  (Behauptung, Vermutung). 

\item Fehler 1. Art ($\alpha$-Fehler) und 2. Art ($\beta$-Fehler)
  
\item Signifikanzniveau\index{Signifikanzniveau} typischerweise 99\%
  o.\,ä. bei binomialverteilten Zufallsvariablen.

\item $p$-Wert
  
\end{itemize}

\newpage

\subsection{Würfelexperiment}

%%
%Bearbeiten Sie mit der Klasse das Würfelexperiment im OLAT:

%\weblink{Würfelexperiment}{https://olat.bms-w.ch/auth/RepositoryEntry/6029794/CourseNode/107682584460080}



\subsubsection{\epsdice{3} werfen}
\TRAINER{Die Würfel werden vorab so sortiert, dass die Hälfte der
Klasse einen gezinkten Würfel mit doppelter \epsdice{3} erhalten wird.}

Werfen Sie den Würfel 29 Mal und notieren Sie das Ergebnis wie folgt:
\begin{itemize}
\item Wie oft kommt die \epsdice{3},
\item wie oft mehr als die \epsdice{2}?
\end{itemize}

\begin{bbwFillInTabular}{|c|p{50mm}|}\hline
$\epsdice{3}$ & \\\hline
$\epsdice{3},\epsdice{4},\epsdice{5}$ oder $\epsdice{6}$ & \\\hline
\epsdice{1} oder \epsdice{2} & \\\hline
\end{bbwFillInTabular}

\subsubsection{Ergebnis}

Wie oft ist die \epsdice{3} erschienen? ....................
\vspace{15mm}

Wie oft kam mehr als die \epsdice{2}? .......................

\newpage


\subsection{Nullhypothese}
Angenommen, der Würfel ist «fair», so erwarten wir, dass jede Seite
etwa gleich häufig auftritt. Hierzu ein Begriff:

\begin{center}\textbf{Nullhypothese} oder {$H_0$}.\end{center}

\begin{definition}{Nullhypothese\index{Nullhypothese}}{}
Die \textbf{Nullhypothese $H_0$} nimmt an, dass bloß der reine Zufall
im Spiel war und sonst nichts.
\end{definition}

\begin{definition}{Alternativhypothese\index{Alternativhypothese}, Forschungshypothese\index{Forschungshypothese}}{}
Die \textbf{Alternativhypothese $H_1$} hingegen behauptet, dass das
Ergebnis nicht allein durch den Zufall zustande kommt; es steht eine
«\textbf{Vermutung}» oder eine «\textbf{Behauptung}» im Raum. 
\end{definition}



\subsubsection*{Aufstellen der Hypothese}

\begin{center} $H_0$ (Nullhypothese):

\TNT{2}{\textbf{Der Würfel ist fair}: Er liefert mit $\frac{1}{6}$
  aller Würfe eine \epsdice{3}. (Bei 29-maligem Werfen
typischerweise nicht mehr als 4 bis 5 Mal die \epsdice{3}.)}%% end TNT
\end{center}

Behauptung:

\begin{center} $H_1$ (Alternativhypothese):

\TNT{2}{\textbf{Der Würfel ist nicht fair}: Er zeigt mehr
  \epsdice{3}-er als der bloße Zufall. Er liefert bei 29-maligem Werfen
mehr als 5 Mal die \epsdice{3}.}
\end{center}


\begin{bemerkung}{Einseitiger Hypothesentest}{}
  Wir sprechen hier von einer \textbf{einseitigen} Hypothese
/ bzw. vom einseitigen Hypothesentest. Ein \textbf{zweiseitiger}
Hypothesentest würde auch das Minimum festlegen.
\end{bemerkung}
\newpage


\newpage
\subsubsection{Erwartete Resultate}

Wie oft erwarten wir nun vom «fairen» Würfel, das Auftreten
einer \epsdice{3} bei 29 maligem Würfeln?

\TNT{4}{
$\frac{29}{6}\approx 4.83$ Mal. Es ist jedoch \textbf{nicht} zu erwarten, dass
auch ein fairer Würfel \textbf{immer} genau 4 oder 5 Mal die \epsdice{3} zeigen wird!
}


Ab wann werden wir nicht mehr glauben, dass der Würfel fair ist?

Zeigt der Würfel ein Resultat, das extrem unwahrscheinlich ist, dann
können wir annehmen, dass der Würfel irgendwie gezinkt ist.

Auch wenn der Würfel nun nicht exakt fünf mal die \epsdice{3} bringt, können
wir eventuell annehmen, dass der Würfel fair ist und ein zweimaliges
Auftreten kann genauso gut eine Laune des Zufalls sein, wie ein
sechsmaliges Auftreten.

Berechnen Sie, wie groß die Wahrscheinlichkeit ist, dass ein «fairer»
Würfel bei 29-maligem Werfen \textbf{maximal} sechs Mal eine \epsdice{3}
zeigt:

\TNT{1.6}{Bernoulli $P(X\le 6 \text{ Mal die }\epsdice{3})  = \sum_{k=0}^6{29  \choose
k} \left(\frac16\right)^k \cdot{} \left(\frac56\right)^{29-k}\approx 80.22\% $}
\newpage
Berechnen Sie mit dem Taschenrechner alle Wahrscheinlichkeiten, dass
$k$ mal eine \epsdice{3} geworfen wird innerhalb von $29$ Würfen.

\begin{bbwFillInTabular}{|c|p{10cm}|}\hline
$k$  & Wahrscheinlichkeit, \textbf{maximal} $k$-mal eine \epsdice{3} zu erzielen
bei 29-maligem werfen: \\\hline\hline
 3   & \TRAINER{26.42 \%}\\\hline
 4   & \TRAINER{45.64 \%}\\\hline
 5   & \TRAINER{64.85 \%}\\\hline
 6   & \TRAINER{80.22 \%}\\\hline
 7   & \TRAINER{90.31 \%}\\\hline
 8   & \TRAINER{95.87 \%}\\\hline
 9   & \TRAINER{98.46 \%}\\\hline
\end{bbwFillInTabular}
\newpage
\subsubsection*{Erwartete Verteilungen}
\bbwCenterGraphic{12cm}{geso/stoch/img/wuerfeltabelle.png}
\newpage

Mit 95.87 \% Wahrscheinlichkeit liefert der normale Würfel bis max. 8 Mal die \epsdice{3}.

Mit 93.16 \% Wahrscheinlichkeit liefert unser gezinkte  Würfel bis
max. 13 Mal die \epsdice{3}. \TRAINER{Diese Berechnung funktioniert
  bei uns natürlich nur, weil die die Art des «Zinkens» kennen.}

\bbwCenterGraphic{17cm}{geso/stoch/img/wuerfelverteilung.png}
\newpage
\subsubsection{Unsere Resultate}

Füllen wir in der Klasse gemeinsam die folgende Tabelle aus:

\begin{bbwFillInTabular}{p{5cm}|p{5cm}|p{5cm}}
    & Würfel ist «fair» ($H_0$) & Würfel ist gezinkt ($H_1$)\\\hline
Würfel wird als «fair» gewertet (weniger als acht \epsdice{3}-er).&
& \TRAINER{Fehler 2. Art}\\\hline
Würfel wird als gezinkt gewertet (mehr als
sieben \epsdice{3}-er).& \TRAINER{Fehler erster Art}& \\\hline
\end{bbwFillInTabular}

\newpage


\subsection{Fehler erster und Fehler zweiter Art}
Wo «kippt» die Erwartung nun wohl am ehesten von einem «fairen» Würfel
zu einem gezinkten, der auf zwei Seiten eine \epsdice{3} stehen hat?

Wo sind wir am «unsichersten»?
Beim etwa \LoesungsRaum{6-8} maligem Auftreten der \epsdice{3} sind wir
wohl am unsichersten.

Legen wir fest: Ein Würfel, der weniger als acht Mal
eine \epsdice{3} zeigt bei 29-maligem würfeln, gelte noch als «fair».

\subsubsection{Signifikanzniveau}
Das Signifikanzniveau wird in Sozialwissenschaften oft als 95 \%
Wahrscheinlichkeit «sich nicht zu irren» angenommen.

Da wir die «Art des Zinkens» bei unseren Würfeln kennen, legen wir das
Signifikanzniveau da fest, wo wir die größte Unsicherheit haben, ob
es ein fairer Würfel mit zufälligerweise vielen \epsdice{3}-ern ist,
oder ob es ein gezinkter Würfel mit ausnahmsweise wenigen
\epsdice{3}-ern ist.

Das Signifikanzniveau legen wir als bei «weniger als acht Mal die
\epsdice{3} bei 29maligem Würfeln» fest (siehe obige Graphik):
\vspace{15mm}
\textbf{Signifikanzniveau} = \LoesungsRaum{90.31\%}


\subsubsection{Fehler erster Art}
Fehler erster Art (Alpha ($\alpha$)-Fehler) = Fälschliche Zurückweisung der Nullhypothese.

Der «faire» Würfel wird als gezinkt kategorisiert.

Dies kann auch in einer Beobachtungsstudie passieren, wenn \zB
angenommen wird, dass die Winterthurer häufiger krank sind als der
Rest der Schweiz, nur weil die Winterthurer in der Statistik
zufälligerweise häufiger Krankheitssymptome zeigen als andere. Dies
kann verschiedene Ursachen haben: Zu kleine Stichprobe und
zufälligerweise einige Kranke erwischt; oder Bias in der Stichprobe,
weil die Umfrage zu nahe am Kantonsspital stattgefunden hat oder ganz
einfach weil das Resultat zufälligerweise zustande kommt etc.

\subsubsection{Fehler zweiter Art}
Fehler zweiter Art (Beta ($\beta$)-Fehler) = Fälschliche Annahme der Alternativhypothese.

Der gezinkte Würfel wird als «fair» angesehen.

\TNTeop{}
\newpage


\subsection{Entscheidungsregel}
\subsubsection*{Annahmebereich / Ablehnungsbreich}\index{Annahmebereich}\index{Ablehnungsbereich}
(Wir betrachten nur Binomialverteilungen)

Wir beobachten, dass der Würfel mehr als angenommen eine \epsdice{3}
zeigt und wir \textbf{vermuten}, dass der Würfel gezinkt ist.

\textbf{Nullhypothese} $H_0$: Der Würfel ist fair und zeigt in
$\frac16$ der Fälle eine \epsdice{3}.

\textbf{Alternativhypothese} $H_1$: Der Würfel zeigt in mehr als
$\frac{1}{6}$ der Fälle eine \epsdice{3}.

Wir wollen für unsere Annahme ($H_1$) maximal mit 5 \%-iger
Wahrscheinlichkeit falsch liegen.

Wir werfen also 29 Mal und wir wollen herausfinden, bis wie viele
\epsdice{3}-er denn noch \textbf{normal} sind:
Also: Mit mit 95 \%-iger Wahrscheinlichkeit kommen 0 bis k
\epsdice{3}-er vor in einem normalen Würfel:

\texttt{binomial\textbf{c}df} $n=29; p=\frac16$

Dies liefert: 0-8 mal kommt die \epsdice{3} mit 95.87 \%-iger
Wahrscheinlichkeit vor. Nur die restlichen ca 4.13\% der
Versuchsreihen kommt die \epsdice{3} mehr als 8 mal vor.


Annahmebereich:

\TNT{2}{0 bis 8 \epsdice{3}-ern und der}

Ablehungsbereich:

\TNT{2}{\textbf{Ablehnungsbereich} liegt bei 9 bis 29 \epsdice{3}-ern.
Wir sprechen hier von einem «rechtsseitigen» Hypothesentest.
}

\newpage


\subsection*{Referenzaufgabe}
Eine groß angelegte Statistik zeigt, dass im Vorschulalter 11 \% der
Mädchen Linkshänderinnen sind. Es wird \textbf{vermutet}, dass nach
vollendeter Schulzeit prozentual weniger Abgängerinnen
Linkshänderinnen sind. Sie haben 45 Schulabgängerinnen zur Überprüfung
Ihrer Hypothese zur Verfügung.

a) Durch einen Hypothesentest will man zeigen, dass bei den
Schulabgängerinnen $P < 0.11$ gilt. Wie lautet die $H_0$-Hypothese
(Nullhypothese)? Und wie die Alternativhypothese ($H_1$), also die
Hypothese zur \textbf{Vermutung}? Gehen Sie von einer
Binomialverteilung aus.

\TNT{4}{
$H_0$ Hypothese: (weiterhin) 11 \% der (Schul-)Abgängerinnen sind
  Linkshänderinnen.

$H_1$ Hypothese: Es sind weniger als 11 \% der (Schul-)Abgängerinnen Linkshänderinnen.
}%% end TNT
 
b) Geben Sie den Annahmebereich von $H_0$ ($[k; 45]$ und den
Ablehnungsbereich von $H_0$ ($[0; k-1]$) der
Nullhypothese auf einem  5\% Signifikanzniveau (Heißt $\alpha \le
5\%$ «Wir können der Entscheidungsregel zu 95 \% vertrauen.») an bei einer Stichprobe von 45 Abgängerinnen.

\TNTeop{
Der Erwartungswert liegt bei 4.95 linkshändigen Mädchen. Wir machen
einen \textbf{linksseitigen} Hypothesentest

Ansatz: $P(X \le k-1) \le 0.05$
  
Binomial\textbf{c}df liefert bei $n\le 1$ noch 3.46 \% und bei $n\le 2$ bereits
11.45 \%. Somit ist der Annahmebereich der $H_0$-Hypothese (also der
Nullhypothese) 2 bis (und mit) 45: $k=2$.

Der Ablehnungsbereich der $H_0$ Hypothese entspricht somit 0 und
1. Wenn also keine oder nur eine Abgängerin Linkshänderin ist, so
können wir mit über 95 \%-iger Wahrscheinlichkeit davon ausgehen, dass
die Nullhypothese verworfen wird. In diesem Fall können wir sagen,
dass mit großer Wahrscheinlichkeit die «Linkshändigkeit» während der
Schulzeit abnimmt.
}

\olatLinkVideo{Hypothesentest}{https://olat.bms-w.ch/auth/RepositoryEntry/6029794/CourseNode/113436965452103}{Detaillierte Erklärung}
%%%%%%%%%%%%%%%%%%%%%%%%%%%%%%%%%%%%%%%%%%%%%%%%%%%%%%%%%%%%%%%%%%%%%%%%5
